\chapter{RAG-Generierung}
\label{ch:rag}
% Deckt Ziel 4 ab: "Auswahl eines fuer die Hardware geeigneten Sprachmodells"
% + die eigentliche Antwortgenerierung.

Nachdem der Retriever aus Kapitel~\ref{ch:indexierung} zu einer Anfrage die am besten
passenden Text-Chunks liefert, müssen diese in eine für ein Sprachmodell verwertbare
Form gebracht und anschließend an ein tatsächlich auf der Zielhardware lauffähiges
Sprachmodell übergeben werden. Dieses Kapitel behandelt die dafür zuständige
Pipeline-Stufe: die Konstruktion des Prompts aus Anfrage und Kontext, die Auswahl
eines für den Raspberry Pi geeigneten Sprachmodells sowie die Anbindung der beiden
unterstützten Inferenz-Backends. Den Abschluss bilden konkrete Beispielinteraktionen,
die das resultierende Grounding-Verhalten des Systems veranschaulichen.

\section{Prompt-Konstruktion}
\label{sec:prompt-konstruktion}

Die Übersetzung von Anfrage und erhaltenen Chunks in einen an das Sprachmodell
sendbaren Prompt übernimmt das Modul \texttt{rag/prompt\_builder.py}. Es enthält
bewusst weder Modell-Ladelogik noch Netzwerk-I/O und lässt sich dadurch unabhängig
von einem laufenden Backend testen. Das Ergebnis ist keine einzelne Zeichenkette,
sondern eine Liste von Chat-Nachrichten der Form
\texttt{[\{"role": ..., "content": ...\}]}, da sowohl die \texttt{llama-cpp-python}-
als auch die Ollama-Schnittstelle (Abschnitt~\ref{sec:backend-anbindung}) dieses
Format erwarten und der Prompt-Builder dadurch backend-unabhängig bleibt.

\paragraph{System-Prompt.} Jeder Anfrage wird derselbe System-Prompt vorangestellt:

\begin{quote}
\ttfamily\small
"You are an assistant answering questions using only the Wikipedia excerpts given
below as context. Answer using only information contained in the excerpts. If the
excerpts do not contain enough information to answer, say so explicitly instead of
guessing. Keep answers concise, and refer to excerpts by their [n] number when
useful."
\end{quote}

Der System-Prompt verfolgt damit zwei Ziele: Erstens soll das Modell seine Antwort
ausschließlich auf den mitgelieferten Wikipedia-Kontext stützen (Grounding) und nicht
auf im Modell selbst gespeichertes Weltwissen zurückgreifen. Zweitens wird das Modell
explizit angewiesen, bei unzureichendem Kontext dies offen zuzugeben, anstatt eine
Antwort zu erraten. Beide Anweisungen lassen sich nicht technisch erzwingen, sondern
sind Instruktionen an ein instruction-getuntes Sprachmodell; ihre Wirksamkeit hängt
entsprechend vom gewählten Modell ab (Abschnitt~\ref{sec:slm-auswahl}) und wird in
Abschnitt~\ref{sec:beispielinteraktionen} anhand konkreter Beispiele eingeordnet.

\paragraph{Kontextformatierung und Quellenattribution.} Die erhaltenen Chunks
werden über \texttt{format\_context()} zu einem nummerierten, quellenattributierten
Textblock zusammengesetzt:

\begin{quote}
\ttfamily\small
[1] (Climate change - Causes) Human activity is the main driver...\\
{[2]} (Climate change - Lead) Present-day climate change includes...
\end{quote}

Jede Zeile enthält neben dem eigentlichen Chunk-Text den Titel des zugehörigen
Wikipedia-Artikels sowie den Titel des Abschnitts, aus dem der Chunk stammt. Beide
Angaben stehen bereits als Metadaten im \texttt{Chunk}-Objekt zur Verfügung, da sie
während des Chunkings pro Abschnitt gesetzt werden (Abschnitt~\ref{sec:chunking-strategie}).
Die fortlaufende Nummerierung erlaubt es dem Sprachmodell, gemäß der Anweisung im
System-Prompt einzelne Aussagen im Fließtext auf eine konkrete Quelle zurückzuführen
(\texttt{[n]}). Dieser Mechanismus macht die Herkunft einer generierten Aussage
nachvollziehbar, ersetzt jedoch keine Prüfung, ob das Modell die Quellenangabe auch
korrekt zuordnet.

\paragraph{Kontext-Budget.} Da erhaltenen Chunks in der Summe das Kontextfenster des
Sprachmodells überschreiten können, begrenzt \texttt{\_fit\_to\_budget()} die
tatsächlich in den Prompt aufgenommene Chunk-Menge auf ein Wortbudget von
standardmäßig 800 Wörtern (\texttt{DEFAULT\_MAX\_CONTEXT\_WORDS}). Die Wortanzahl dient
dabei als Näherung für die tatsächliche Token-Anzahl: Für englischsprachigen Text
liegen gängige Tokenizer bei ungefähr 1,3 Token pro Wort, sodass 800 Wörter deutlich
unterhalb des in \texttt{config.llm.n\_ctx} konfigurierten Kontextfensters von 4096
Token bleiben. Der verbleibende Spielraum wird für den System-Prompt sowie die
maximale Antwortlänge (\texttt{config.llm.max\_tokens}, standardmäßig 400 Token)
benötigt.

Die Auswahl der tatsächlich in den Prompt aufgenommenen Chunks erfolgt greedy: Die
bereits nach Relevanz sortierten Chunks werden in dieser Reihenfolge durchlaufen,
wobei ein Chunk, der das Budget überschreiten würde, übersprungen und die Prüfung mit
dem nächstplatzierten Chunk fortgesetzt wird, statt beim ersten nicht mehr passenden
Chunk abzubrechen. Dasselbe Prinzip verwendet bereits die budgetbeschränkte
Artikelauswahl aus Kapitel~\ref{ch:dokumentenauswahl}; ein einzelner, ungewöhnlich
langer Chunk verdrängt dadurch nicht automatisch alle nachfolgenden, kürzeren und
unter Umständen ebenfalls relevanten Chunks.

\paragraph{Leere Retrieval-Ergebnisse.} Liefert der Retriever keine Chunks -- etwa bei
einer Anfrage zu einem nicht im Index enthaltenen Thema --, wird der Prompt dennoch
abgesendet. An die Stelle des Kontextblocks tritt dann der feste Hinweistext
\texttt{No relevant excerpts were found in the indexed Wikipedia subset.} Ein
Sonderfall im Aufrufercode ist dafür nicht nötig, da der System-Prompt das Modell
bereits anweist, in einer solchen Situation das Fehlen ausreichender Informationen zu
benennen. Wie sich dieses Zusammenspiel in der Praxis auswirkt, wenn zwar keine
vollständig leere, aber eine inhaltlich unpassende Trefferliste zurückgegeben wird,
wird in Abschnitt~\ref{sec:beispielinteraktionen} diskutiert.


\section{Sprachmodell-Auswahl}
\label{sec:slm-auswahl}

Für die Textgenerierung wurden drei quantisierte, instruction-getunte Sprachmodelle
im GGUF-Format evaluiert: \texttt{Qwen2.5-1.5B-Instruct}~\cite{qwen25-model-card},
\texttt{Llama-3.2-1B-Instruct}~\cite{llama32-model-card} und
\texttt{Phi-3.5-mini-instruct}~\cite{phi35-model-card}. Alle drei Kandidaten liegen im
Q4\_K\_M-Quantisierungsformat vor, einer 4-Bit-Gewichtsquantisierung, die von der für
die CPU-only-Inferenz auf dem Raspberry Pi eingesetzten
\texttt{llama.cpp}-Bibliothek~\cite{llamacpp-repo} unterstützt wird und gegenüber der
vollpräzisen Modellversion Speicherbedarf und Rechenaufwand deutlich reduziert. Die
Auswahl dieser drei Modelle war durch die Zielhardware motiviert: Alle drei zählen zu
den kleineren aktuell verfügbaren instruction-getunten Sprachmodellen und werden von
ihren jeweiligen Herausgebern explizit auch für ressourcenbeschränkte Einsatzszenarien
vorgesehen.

\begin{table}[ht]
    \centering
    \caption{Verglichene Sprachmodelle}
    \label{tab:slm-vergleich}
    \begin{tabular}{lrrrl}
        \toprule
        Modell & Parameter & Kontextfenster & Größe (Q4\_K\_M) & Lizenz \\
        \midrule
        Llama 3.2 1B Instruct   & 1,23 Mrd. & 128.000 Token & 771 MB & Llama 3.2 Community License \\
        Qwen2.5 1.5B Instruct   & 1,54 Mrd. & 32.768 Token  & 1,1 GB & Apache 2.0 \\
        Phi-3.5 Mini Instruct   & 3,8 Mrd.  & 128.000 Token & 2,3 GB & MIT \\
        \bottomrule
    \end{tabular}
\end{table}

Die in Tabelle~\ref{tab:slm-vergleich} angegebenen Kontextfenster sind die von den
Herausgebern dokumentierten nativen Maximalwerte der jeweiligen Modelle. Innerhalb von
WikiPod wird davon unabhängig für alle drei Modelle einheitlich das in
\texttt{config.llm.n\_ctx} konfigurierte Kontextfenster von 4096 Token verwendet
(Abschnitt~\ref{sec:prompt-konstruktion}). Ein größeres Kontextfenster würde bei
\texttt{llama.cpp} zusätzlichen Arbeitsspeicher für den Attention-Cache (KV-Cache)
belegen, ohne dass die Prompt-Konstruktion aus Abschnitt~\ref{sec:prompt-konstruktion}
-- System-Prompt, ein auf 800 Wörter begrenzter Kontextblock und eine auf 400 Token
begrenzte Antwort -- diesen Spielraum überhaupt ausschöpfen würde. Auf der
CPU-only-Zielhardware wird das konfigurierte Kontextfenster somit bewusst klein
gehalten.

Die eigentliche Generierungslatenz der drei Modelle auf dem Raspberry Pi wurde bereits
in Kapitel~\ref{ch:evaluation} (Tabelle~\ref{tab:pi-generation-latency}) im Rahmen
der Effizienzevaluation gemessen und dort ausführlich diskutiert: Qwen2.5 1.5B und
Llama 3.2 1B lagen mit rund 28 beziehungsweise 29 Sekunden mittlerer Latenz pro
Anfrage nah beieinander, während Phi-3.5 Mini mit über 150 Sekunden deutlich
langsamer war. Dieser Abschnitt konzentriert sich daher ergänzend auf die
Antwortqualität, da für den in Kapitel~\ref{ch:evaluation} dokumentierten
Modellvergleich bewusst keine quantitative Qualitätsmetrik erhoben wurde.

Bei manueller Durchsicht der bei der Latenzmessung erzeugten Antworten desselben
Evaluationsdatensatzes zeigen sich zwischen den drei Modellen qualitative
Unterschiede im Umgang mit dem im System-Prompt geforderten Antwortstil. Auf die
Anfrage \textit{"Which music group did Michael Jackson perform with before becoming a
successful solo artist?"} antwortete Qwen2.5 1.5B mit \texttt{The Jackson 5} -- knapp
und ohne überflüssige Ausführungen, wie es der System-Prompt (\texttt{Keep answers
concise}) vorgibt. Phi-3.5 Mini beantwortete dieselbe Anfrage inhaltlich korrekt,
jedoch deutlich ausführlicher und mit explizit in den Fließtext eingebundenen
Quellenverweisen (\textit{"According to the information provided in excerpts [1] and
[2], Michael Jackson performed with The Jackson 5 [...]"}). Llama 3.2 1B lieferte für
dieselbe Anfrage eine inhaltlich ebenfalls korrekte Antwort, die jedoch mit der
unverändert aus dem Kontextblock übernommenen Nummerierung \texttt{[1]} beginnt
(\textit{"[1] The Jackson 5, later known as The Jacksons, was an American pop band
[...]"}), anstatt den Quellenverweis wie vom System-Prompt vorgesehen in den
Antworttext zu integrieren. Ein vergleichbares Muster zeigt sich bei der Anfrage nach
dem südafrikanischen Anti-Apartheid-Aktivisten, der 27 Jahre im Gefängnis saß, bevor
er ins Präsidentenamt gewählt wurde: Llama 3.2 1B beantwortete auch diese Anfrage mit
vorangestelltem \texttt{[1]}, während Qwen2.5 1.5B inhaltlich nahezu identisch, aber
ohne diesen Formatierungsartefakt antwortete.

Diese Beobachtungen beruhen auf einer kleinen, nicht systematisch ausgewerteten
Stichprobe einzelner Antworten und ersetzen keine quantitative Qualitätsbewertung;
sie werden hier ausschließlich als Begründung für die getroffene Auswahlentscheidung
angeführt. In der Zusammenschau aus gemessener Latenz und den beschriebenen
qualitativen Beobachtungen wurde \texttt{Qwen2.5-1.5B-Instruct} als Standardmodell für
\texttt{config/default.yaml} und \texttt{config/prod.yaml} übernommen: Es erreichte in
Kapitel~\ref{ch:evaluation} die niedrigste gemessene Generierungslatenz der drei
Kandidaten, produzierte in der hier beschriebenen Stichprobe durchgehend prägnante,
inhaltlich mit den beiden anderen Modellen übereinstimmende Antworten und wies dabei
nicht das bei Llama 3.2 1B beobachtete Formatierungsartefakt auf. Phi-3.5 Mini wurde
trotz ausführlicherer, explizit belegter Antworten aufgrund der deutlich höheren
Latenz nicht als Standardmodell übernommen.

\todo{Vollständiger End-to-End-Lauf auf dem Edge-Case-Datensatz
für alle drei Modelle bauen und zeige und auch die Grenzfälle aus Kapitel~\ref{ch:evaluation}}


\section{Backend-Anbindung}
\label{sec:backend-anbindung}

Die Anbindung des Sprachmodells kapselt die Klasse \texttt{Generator} in
\texttt{rag/generator.py}. Sie nimmt die von \texttt{build\_messages()} erzeugte
Chat-Nachrichtenliste entgegen und gibt unabhängig vom konfigurierten Backend einen
generierten Antworttext zurück, sodass der aufrufende Code (\texttt{cli.py}) nicht
zwischen den Backends unterscheiden muss. Welches Backend verwendet wird, legt
\texttt{config.llm.backend} fest; unterstützt werden zwei Werte:

\begin{itemize}
    \item \texttt{'llama\_cpp'}: Lädt eine lokale GGUF-Datei direkt über die
    \texttt{llama-cpp-python}-Bindung (eine optionale Abhängigkeit, siehe das
    \texttt{[llm]}-Extra in \texttt{pyproject.toml}) und ruft
    \texttt{create\_chat\_completion()} mit den konfigurierten Werten für
    \texttt{max\_tokens} und \texttt{temperature} auf. Dieses Backend läuft ohne
    separaten Server-Prozess und entspricht damit dem für die Zielhardware
    vorgesehenen Deployment: In \texttt{config/prod.yaml} ist ausschließlich dieses
    Backend konfiguriert.
    \item \texttt{'ollama'}: Spricht stattdessen per HTTP einen laufenden
    Ollama-Daemon an (\texttt{POST /api/chat}) und übergibt Nachrichtenliste sowie
    Generierungsparameter als JSON-Body. Da für dieses Backend lediglich
    \texttt{ollama pull <modell>} statt eines manuellen GGUF-Downloads und Kompilierens
    von \texttt{llama-cpp-python} nötig ist, wird es in \texttt{config/dev.yaml} und
    \texttt{config/pi-test.yaml} für die lokale Entwicklung beziehungsweise einen
    Testlauf im mittleren Maßstab verwendet.
\end{itemize}

Beide Implementierungen erwarten und liefern dasselbe von
Abschnitt~\ref{sec:prompt-konstruktion} beschriebene Nachrichtenformat, sodass ein
Wechsel des Backends -- etwa von \texttt{ollama} während der Entwicklung zu
\texttt{llama\_cpp} auf dem Raspberry Pi -- keine Änderung an Prompt-Konstruktion oder
Retrieval erfordert. Das Laden eines GGUF-Modells über \texttt{llama\_cpp} ist der mit
Abstand teuerste Teil der Backend-Initialisierung; \texttt{\_load\_llama\_cpp\_model()}
cached daher analog zum in Kapitel~\ref{ch:indexierung} beschriebenen Embedding-Modell
die geladene \texttt{Llama}-Instanz pro Kombination aus Modellpfad und Kontextgröße
(\texttt{@lru\_cache(maxsize=2)}), sodass ein Modell innerhalb eines Prozesses nur
einmal geladen werden muss. Fehlt die konfigurierte GGUF-Datei, bricht der
\texttt{llama\_cpp}-Pfad mit einer aussagekräftigen Fehlermeldung ab, die sowohl auf
eine Bezugsquelle für das Modell als auch auf die Möglichkeit verweist, für die lokale
Entwicklung stattdessen auf das \texttt{ollama}-Backend zu wechseln.


\section{Beispielinteraktionen}
\label{sec:beispielinteraktionen}

Anhand einer Anfrage aus dem regulären Evaluationsdatensatz
(Kapitel~\ref{ch:evaluation}) lässt sich das Grounding-Verhalten des Systems
konkret nachvollziehen. Auf die Anfrage \textit{"Which South African anti-apartheid
activist spent 27 years in prison before becoming the country's first black
president?"} lieferte das mit \texttt{Qwen2.5-1.5B-Instruct} generierte System die
Antwort:

\begin{quote}
\itshape
Nelson Rolihlahla Mandela spent 27 years in prison before becoming the country's
first black president.
\end{quote}

Die Antwort entspricht dabei nahezu wörtlich einer Formulierung aus dem extrahierten
Wikipedia-Artikel zu Nelson Mandela, dem gemäß Evaluationsdatensatz auch tatsächlich
relevanten Ground-Truth-Artikel. Sie belegt damit, dass das Modell die Antwort aus dem
im Prompt bereitgestellten Kontext übernimmt, statt auf einen aus dem Vortraining
stammenden, nicht mit einer Quelle belegten Wissensstand zurückzugreifen.

Für den in der Aufgabenstellung ebenfalls geforderten negativen Fall -- eine Anfrage zu
einem nicht im indexierten Bestand enthaltenen Thema -- enthält der in
Kapitel~\ref{ch:evaluation} beschriebene Edge-Case-Datensatz
(\texttt{test/data/eval\_edge\_case\_queries.yaml}) die als \texttt{out\_of\_scope}
markierte Anfrage \texttt{Tell me about Albert Einstein}, für die keine relevanten
Ground-Truth-Titel hinterlegt sind. Die dort dokumentierte Retrieval-Auswertung zeigt,
dass für diese Anfrage keine leere Trefferliste zurückgegeben wird, sondern
inhaltlich nicht einschlägige Chunks -- unter anderem zum Artikel \texttt{Martin
Luther King Jr.} Für diesen konkreten Fall greift somit nicht die in
Abschnitt~\ref{sec:prompt-konstruktion} beschriebene Fallback-Behandlung für leere
Retrieval-Ergebnisse, sondern es liegt am System-Prompt selbst, anhand eines
inhaltlich unpassenden, aber nicht leeren Kontextblocks zu erkennen, dass die
Anfrage nicht ausreichend beantwortet werden kann.

